\section{Conclusion}\label{sec:conclusion}
We asked when a phishing detector facing LLM-generated mail should buy its scarce labels. A label-free sliced test on a frozen representation, with a level schedule that bounds false alarms over a whole reference epoch, answers this well. On a 36{,}000-message stream it reached MCC {{main|SLIDE|mcc:.3f}} with {{main|SLIDE|labels:,.0f}} labels, {{=[main|SLIDE|mcc]-[main|Budget-Uniform-MLP|mcc]:.3f}}--{{=[main|SLIDE|mcc]-[main|Budget-Random|mcc]:.3f}} above equal-budget samplers that ignore the shift and within {{=abs([vs|Budget-Oracle|d]):.3f}} of an oracle. The same ordering held for a TF-IDF learner and four stream designs, and the trigger raised no false alarm on 20 stationary streams. A fully supervised online TF-IDF learner remains more accurate when labels are plentiful and arrive quickly. The trigger is blind to posterior change with unchanged inputs, which a labelled audit of a tenth of the traffic only partly repaired, and the 128-token encoder is defeated by benign padding unless messages are scored in windows at a cost in false positives. Training the head on all windows, combining the trigger with a whole-message learner, and testing it on traffic recorded at a single organisation are the next steps.
